En esta práctica se estudió dónde se forma la imagen de un objeto frente a un espejo plano. Sobre un pliego de papel se colocó un espejo de pie a lo largo de una línea central. Se pusieron objetos en once posiciones distintas frente a él, y para cada una se localizó su imagen detrás del espejo con el método de paralaje, midiendo después la distancia del objeto al espejo y la de la imagen al espejo. La gráfica de una distancia contra la otra dio una recta con pendiente de 1.005 y un coeficiente de determinación de 0.9996. En promedio la imagen quedó apenas 0.24 cm más cerca del espejo que el objeto, con un error relativo medio de 1.8 por ciento. Esto confirma que en un espejo plano la imagen está a la misma distancia detrás del espejo que el objeto delante de él.
